\subsubsection{Functional Requirements}

\resetreqcounter{GeFR}
\begin{srstable}[
  caption = {Общие функциональные требования к Сайту.},
]{colspec = {Q[c, m] X[j, m]}}
  Идентификатор & Описание \\
  \reqlabel[ref:gef1]{GeFR} & Сайт \textit{Должен} предоставлять следующие \mbox{элементы} в Шапке:
  \begin{compactitem}
    \item логотип Сайта, ведущий на главную страницу;
    \item список доступных гаданий, переводящих на соответсвующие страницы гаданий;
    \item поиск по содержимому страниц Сайта.
  \end{compactitem} \\
  \reqlabel[ref:gef2]{GeFR} & Сайт \textit{Должен} показывать Пользователям рекламу следующих форматов:
  \begin{compactitem}
    \item \textit{баннерная:} между блоками Сайта;
    \item \textit{всплывающая:} в углу экрана.
  \end{compactitem} \\
  \reqlabel[ref:gef3]{GeFR} & Сайт \textit{Должен} вести статистику посещений Пользователей. \\
\end{srstable}

\resetreqcounter{MpFR}
\begin{srstable}[
  caption = {Функциональные требования к главной странице Сайта.},
]{colspec = {Q[c, m] X[j, m]}}
  Идентификатор & Описание \\
  \reqlabel[ref:mpf1]{MpFR} & Сайт \textit{Должен} содержать блок с новыми способами гадания на главной странице. \\
  \reqlabel[ref:mpf2]{MpFR} & Сайт \textit{Должен} содержать на главной странице раздел со способами гадания, которые соответствуют следующим критериям:
  \begin{compactitem}
    \item календарная дата;
    \item время года;
    \item день недели.
  \end{compactitem} \\
\end{srstable}

\resetreqcounter{LgFR}
\begin{srstable}[
  caption = {Юридические функциональные требования к Сайту.},
]{colspec = {Q[c, m] X[j, m]}}
  Идентификатор & Описание \\
  \reqlabel[ref:lgf1]{LgFR} & Сайт \textit{Должен} содержать страницу с Пользовательским соглашением. \\
  \reqlabel[ref:lgf2]{LgFR} & Сайт \textit{Должен} содержать страницу с Политикой конфиденциальности. \\
  \reqlabel[ref:lgf3]{LgFR} & Сайт \textit{Должен} содержать страницу со своим описанием для Пользователя. \\
  \reqlabel[ref:lgf4]{LgFR} & Сайт \textit{Должен} содержать страницу с контактами Администраторов в случае возникновения вопросов. \\
\end{srstable}

\resetreqcounter{DvFR}
\begin{srstable}[
  caption = {Функциональные требования к процессам гадания на Сайте.},
]{colspec = {Q[c, m] X[j, m]}}
  Идентификатор & Описание \\
  \reqlabel[ref:dvf1]{DvFR} & Сайт \textit{Должен} предоставлять гадания, которые предусматривают один из следующих методов взаимодействия с Пользователем для получения результата:
  \begin{compactitem}
    \item ввод текстового вопроса в свободной форме;
    \item выбор одного или нескольких вариантов из предложенного списка;
    \item нажатие на кнопку без дополнительного ввода;
    \item выбор элемента из визуального ряда.
  \end{compactitem} \\
  \reqlabel[ref:dvf2]{DvFR} & Сайт \textit{Должен} предоставить толкование результата, полученного в процессе гадания. \\
  \reqlabel[ref:dvf3]{DvFR} & Сайт \textit{Должен} сгруппировать гадания следующим образом:
  \begin{compactitem}
    \item \textit{по тематике:} бытовые, любовные и пр.;
    \item \textit{по разделу:} оракул, карты Таро, руны и пр.
  \end{compactitem} \\
  \reqlabel[ref:dvf4]{DvFR} & Сайт \textit{Должен} предоставить следующие критерии сортировки для группы:
  \begin{compactitem}
    \item по дате появления гадания на Сайте;
    \item по популярности гадания на Сайте;
    \item по рейтингу гадания на Сайте.
  \end{compactitem} \\
  \reqlabel[ref:dvf5]{DvFR} & Сайт \textit{Должен} содержать для каждого гадания блоки с описанием, что включает следующее:
  \begin{compactitem}
    \item название услуги;
    \item цель услуги;
    \item формат предоставления услуги;
    \item результат оказания услуги.
  \end{compactitem} \\
\end{srstable}

\resetreqcounter{FbFR}
\begin{srstable}[
  caption = {Функциональные требования к блоку обратной связи.},
]{colspec = {Q[c, m] X[j, m]}}
  Идентификатор & Описание \\
  \reqlabel[ref:fbf1]{FbFR} & Сайт \textit{Должен} оснащать свои страницы блоком обратной связи, который доступен Пользователям. \\
  \reqlabel[ref:fbf2]{FbFR} & Сайт \textit{Должен} предоставлять блок обратной связи, который содержит следующие элементы:
  \begin{compactitem}
    \item рейтинг текущей страницы;
    \item форма отправки комментария;
    \item список опубликованных комментариев;
    \item кнопки для публикации в социальных сетях ссылки на страницу;
  \end{compactitem} \\
  \reqlabel[ref:fbf3]{FbFR} & Сайт \textit{Должен} предоставлять рейтинг, выражаемый в десятибалльной шкале, где $1$ --- неудовлетворительно, $10$ --- замечательно. \\
  \reqlabel[ref:fbf4]{FbFR} & Сайт \textit{Должен} предоставлять Пользователям возможность влиять на рейтинг любой доступной страницы Сайта. \\
  \reqlabel[ref:fbf5]{FbFR} & Сайт \textit{Должен} предоставлять Пользователям возможность оставлять комментарии при помощи формы. \\
  \reqlabel[ref:fbf6]{FbFR} & Сайт \textit{Должен} предоставлять следующие поля для блока комментариев:
  \begin{compactitem}
    \item псевдоним автора комментария;
    \item содержание комментария.
  \end{compactitem} \\
  \reqlabel[ref:fbf7]{FbFR} & Сайт \textit{Должен} предоставлять формы отправки комментариев с полями, которые обязательны к заполнению и проверяются на валидность, то есть на наличие непустых символов. \\
  \reqlabel[ref:fbf8]{FbFR} & Сайт \textit{Должен} предоставлять списком опубликованные комментарии со следующей информацией:
  \begin{compactitem}
    \item дата публикации комментария;
    \item псевдоним автора комментария;
    \item содержание комментария.
  \end{compactitem} \\
  \reqlabel[ref:fbf9]{FbFR} & Сайт \textit{Должен} сортировать опубликованные комментарии по убыванию даты публикации, то есть от новых к старым. \\
  \reqlabel[ref:fbf10]{FbFR} & Сайт \textit{Должен} отображать Пользователю список комментариев с применением Пагинации. \\
\end{srstable}

\resetreqcounter{FvFR}
\begin{srstable}[
  caption = {Функциональные требования к модулю избранных страниц.},
]{colspec = {Q[c, m] X[j, m]}}
  Идентификатор & Описание \\
  \reqlabel[ref:fvf1]{FvFR} & Сайт \textit{Должен} оснащать страницы панелью избранных страниц. \\
  \reqlabel[ref:fvf2]{FvFR} & Сайт \textit{Должен} оснащать панель избранных страниц следующими элементами:
  \begin{compactitem}
    \item кнопка добавления текущей страницы в избранное;
    \item ссылка на страницу «Избранное» с другими избранными страницами.
  \end{compactitem} \\
  \reqlabel[ref:fvf3]{FvFR} & Сайт \textit{Должен} оснащать страницу «Избранное» индивидуальным списком избранных страниц Пользователя. \\
  \reqlabel[ref:fvf4]{FvFR} & Сайт \textit{Должен} оснащать элемент из списка избранных страниц следующими компонентами:
  \begin{compactitem}
    \item название страницы;
    \item описание страницы;
    \item иконка страницы;
    \item кнопка удаления страницы из списка;
    \item кнопка просмотра страницы.
  \end{compactitem} \\
  \reqlabel[ref:fvf5]{FvFR} & Сайт \textit{Должен} страницу «Избранное» содержать блок рекомендаций с гаданиями, имеющими наиболее высокий рейтинг. \\
\end{srstable}

\resetreqcounter{AmFR}
\begin{srstable}[
  caption = {Административные функциональные требования к Сайту.},
]{colspec = {Q[c, m] X[j, m]}}
  Идентификатор & Описание \\
  \reqlabel[ref:amf1]{AmFR} & Сайт \textit{Должен} предоставлять Администратору возможность создавать, просматривать, редактировать и удалять страницы Сайта. \\
  \reqlabel[ref:amf3]{AmFR} & Сайт \textit{Должен} предоставлять Администратору возможность выбрать тип взаимодействия с Пользователем из списка поддерживаемых методов при создании или редактировании способа гадания. \\
  \reqlabel[ref:amf4]{AmFR} & Сайт \textit{Должен} предоставлять Администратору изменять параметры специальных блоков и разделов Сайта. \\
\end{srstable}

\resetreqcounter{NvFR}
\begin{srstable}[
  caption = {Функциональные требования к навигации.},
]{colspec = {Q[c, m] X[j, m]}}
  Идентификатор & Описание \\
  \reqlabel[ref:nvf1]{NvFR} & Сайт \textit{Должен} предоставлять кнопку «Наверх», которая появляется при прокрутке страницы вниз. \\
  \reqlabel[ref:nvf2]{NvFR} & Сайт \textit{Должен} при нажатии кнопки «Наверх» перемещать Пользователя в верхнюю часть страницы. \\
\end{srstable}

\resetreqcounter{AdFR}
\begin{srstable}[
  caption = {Функциональные требования к блоку «Совет для вас».},
]{colspec = {Q[c, m] X[j, m]}}
  Идентификатор & Описание \\
  \reqlabel[ref:adf1]{AdFR} & Сайт \textit{Должен} отображать в нижней части каждой страницы блок «Совет для вас». \\
  \reqlabel[ref:adf2]{AdFR} & Сайт \textit{Должен} оснащать блок «Совет для вас» следующими элементами:
  \begin{compactitem}
    \item название карты Таро;
    \item изображение карты Таро;
    \item текст совета.
  \end{compactitem} \\
  \reqlabel[ref:adf3]{AdFR} & Сайт \textit{Должен} загружать содержимое блока асинхронно при загрузке страницы. \\
  \reqlabel[ref:adf4]{AdFR} & Сайт \textit{Должен} обновлять содержимое блока при каждой перезагрузке страницы, выбираясь случайным образом из пула советов. \\
\end{srstable}

\resetreqcounter{ItFR}
\begin{srstable}[
  caption = {Функциональные требования к блоку «Это интересно».},
]{colspec = {Q[c, m] X[j, m]}}
  Идентификатор & Описание \\
  \reqlabel[ref:itf1]{ItFR} & Сайт \textit{Должен} отображать на страницах способов гаданий блок «Это интересно». \\
  \reqlabel[ref:itf2]{ItFR} & Сайт \textit{Должен} оснащать блок «Это интересно» текстом интересного факта, выбираемого случайным образом из пула фактов. \\
  \reqlabel[ref:itf3]{ItFR} & Сайт \textit{Должен} оснащать блок «Это интересно» на сервере и отображать его при первой загрузке страницы. \\
  \reqlabel[ref:itf4]{ItFR} & Сайт \textit{Должен} оснащать блок «Это интересно» кнопкой для обновления факта без перезагрузки страницы. \\
\end{srstable}

\resetreqcounter{StFR}
\begin{srstable}[
  caption = {Функциональные требования к блоку «Упавшая звезда».},
]{colspec = {Q[c, m] X[j, m]}}
  Идентификатор & Описание \\
  \reqlabel[ref:stf1]{StFR} & Сайт \textit{Должен} с заданной периодичностью отображать на главной странице анимацию «падающей звезды». \\
  \reqlabel[ref:stf2]{StFR} & Сайт \textit{Должен} отображать текст с указанием даты и времени последнего появления звезды. \\
\end{srstable}

\resetreqcounter{RcFR}
\begin{srstable}[
  caption = {Функциональные требования к блоку рекомендаций.},
]{colspec = {Q[c, m] X[j, m]}}
  Идентификатор & Описание \\
  \reqlabel[ref:rcf1]{RcFR} & Сайт \textit{Должен} отображать блок «Судьба и случай рекомендуют» со ссылкой на случайное гадание. \\
  \reqlabel[ref:rcf2]{RcFR} & Сайт \textit{Должен} изменять содержимое блока «Судьба и случай рекомендуют» при каждой перезагрузке страницы, выбирая его случайным образом из пула рекомендаций. \\
\end{srstable}

\resetreqcounter{ArFR}
\begin{srstable}[
  caption = {Функциональные требования к разделу статей.},
]{colspec = {Q[c, m] X[j, m]}}
  Идентификатор & Описание \\
  \reqlabel[ref:arf1]{ArFR} & Сайт \textit{Должен} предоставлять раздел «Статьи». \\
  \reqlabel[ref:arf3]{ArFR} & Сайт \textit{Должен} перенаправлять на страницу статьи при нажатии на ее заголовок. \\
\end{srstable}

\resetreqcounter{FdFR}
\begin{srstable}[
  caption = {Функциональные требования к разделу фонда «Содействие».},
]{colspec = {Q[c, m] X[j, m]}}
  Идентификатор & Описание \\
  \reqlabel[ref:fdf1]{FdFR} & Сайт \textit{Должен} предоставлять раздел «Содействие» с информацией о благотворительном фонде помощи детям. \\
  \reqlabel[ref:fdf3]{FdFR} & Сайт \textit{Должен} оснащать каждую новость заголовком, являющимся внешней ссылкой на основную статью, и аннотацией. \\
\end{srstable}

\resetreqcounter{NmFR}
\begin{srstable}[
  caption = {Функциональные требования к разделу «Значение имени».},
]{colspec = {Q[c, m] X[j, m]}}
  Идентификатор & Описание \\
  \reqlabel[ref:nmf1]{NmFR} & Сайт \textit{Должен} предоставлять раздел «Значение имени». \\
  \reqlabel[ref:nmf2]{NmFR} & Сайт \textit{Должна} оснащать раздел «Значение имени» списком всех доступных имен. \\
  \reqlabel[ref:nmf3]{NmFR} & Сайт \textit{Должен} предоставить возможность фильтрации списка имён по первой букве и полу (мужские/женские). \\
  \reqlabel[ref:nmf4]{NmFR} & Сайт \textit{Должен} предоставлять отдельную страницу для каждого имени. \\
  \reqlabel[ref:nmf6]{NmFR} & Сайт \textit{Должен} предоставлять список имен, подходящих конкретному знаку зодиака. \\
\end{srstable}

\resetreqcounter{NtFR}
\begin{srstable}[
  caption = {Функциональные требования к разделу «Натальная карта».},
]{colspec = {Q[c, m] X[j, m]}}
  Идентификатор & Описание \\
  \reqlabel[ref:ntf1]{NtFR} & Сайт \textit{Должен} предоставлять раздел «Натальная карта». \\
  \reqlabel[ref:ntf2]{NtFR} & Сайт \textit{Должна} предоставлять форму для расчета натальной карты со следующими полями:
  \begin{compactitem}
    \item дата рождения с графическим календарем для выбора;
    \item время рождения с выпадающими списками для часов и минут;
    \item город рождения с автодополнением при вводе первых трех букв;
    \item система расчета с выпадающим списком.
  \end{compactitem} \\
  \reqlabel[ref:ntf3]{NtFR} & Сайт \textit{Должен} асинхронно отображать результат расчета натальной карты после отправки формы. \\
\end{srstable}

\resetreqcounter{MtFR}
\begin{srstable}[
  caption = {Функциональные требования к разделу «Славянская мифология».},
]{colspec = {Q[c, m] X[j, m]}}
  Идентификатор & Описание \\
  \reqlabel[ref:mtf1]{MtFR} & Сайт \textit{Должен} предоставлять раздел «Славянская мифология». \\
  \reqlabel[ref:mtf2]{MtFR} & Сайт \textit{Должен} оснащать раздел «Славянская мифология» информацией о богах, мифологических существах и рунах. \\
  \reqlabel[ref:mtf3]{MtFR} & Сайт \textit{Должен} предоставлять отдельную страницу для каждого мифического объекта. \\
  \reqlabel[ref:mtf4]{MtFR} & Сайт \textit{Должен} оснащать страницу мифического объекта подробной справочной информацией о нем. \\
\end{srstable}

\resetreqcounter{OmFR}
\begin{srstable}[
  caption = {Функциональные требования к разделу «Народные приметы».},
]{colspec = {Q[c, m] X[j, m]}}
  Идентификатор & Описание \\
  \reqlabel[ref:omf1]{OmFR} & Сайт \textit{Должен} предоставлять раздел «Народные приметы» с возможностью выбора конкретной даты. \\
  \reqlabel[ref:omf2]{OmFR} & Сайт \textit{Должен} отображать следующую информацию для выбранной даты:
  \begin{compactitem}
    \item приметы;
    \item именины;
    \item толкование снов;
    \item талисманы.
  \end{compactitem} \\
\end{srstable}